\subsection{What the Transfer Costs to Store}
\label{sec:transfer_cost}
Appendix~\ref{app:error_budget} prices the compression step against the
error that each product already carries. This section prices the transfer
itself. The same rate--distortion protocol is run on carrier tiles rebuilt
from the same native windows by two further transfers: the PSF-aware one,
which restores small-scale structure by construction, and the
inverse-distance average over the four nearest cells that
Section~\ref{sec:resampler} abandoned, which restores none and serves as the
control. The library calls the latter \texttt{linear}, but it is a
Shepard-type weighted average whose kernel width defaults to the
\emph{output} cell scale; Table~\ref{tab:kernel_ablation} finds it $9.5$
times less accurate than Clough--Tocher against an analytic field, while
leaving no excess top-octave power. The control is what makes the question
falsifiable: a deconvolved tile costing more than a Clough--Tocher tile
would establish nothing on its own, since every property of a deconvolving
transfer is confounded with its not being Clough--Tocher. The PSF tiles are
built at the damping at which the reconstruction error was measured
($\lambda=10^{-3}$) and at the nominal \SI{12.5}{m} response width.
Coverage is 40 regions for the PSF transfer and 39 for the inverse-distance
average.

Table~\ref{tab:transfer_cost} reports the three carriers under the short
names used from here on: C--T for Clough--Tocher, PSF for the PSF-aware
transfer, IDW for the inverse-distance average. Every ratio at a given
target is a paired median over the same regions, those bracketed on all
three sweeps there, so that the columns may be read against one another.
The left block gives the compressed size relative to the Clough--Tocher
tile of the same region at the same absolute RMSE, its last column setting
the deconvolving transfer against the non-deconvolving one; the right block
gives each product against the equal-area native crop at
$\gamma_{\mathrm{tile}}=1.3$, composed per region rather than as a ratio of
medians, so that a value below one favours HEALPix.

\begin{table*}[t]
\centering
\caption{What the transfer costs to store}
\label{tab:transfer_cost}
\footnotesize
\begin{tabular}{@{}lrrrrrrr@{}}
\toprule
 & & \multicolumn{3}{c}{Relative to C--T} & \multicolumn{3}{c}{HPX/UTM at $\gamma_{\mathrm{tile}}{=}1.3$} \\
\cmidrule(lr){3-5}\cmidrule(lr){6-8}
Target RMSE & $n$ & PSF & IDW & PSF/IDW & C--T & PSF & IDW \\
\midrule
$0.0001$ & 39 & 1.27 & 1.20 & 1.05 & 1.59 & 2.02 & 1.91 \\
$0.0003$ & 39 & 1.36 & 1.31 & 1.05 & 1.50 & 2.03 & 1.93 \\
$0.001$ & 38 & 1.64 & 1.60 & 1.03 & 1.23 & 2.03 & 1.94 \\
$0.003$ & 36 & 2.08 & 1.94 & 1.06 & 0.93 & 2.01 & 1.90 \\
$0.01$ & 30 & 2.45 & 2.18 & 1.13 & 0.72 & 1.75 & 1.63 \\
\bottomrule
\end{tabular}

\end{table*}


The control is informative (Table~\ref{tab:transfer_cost}). At the same
coding RMSE, each measured against that transfer's own uncompressed
carrier, the deconvolved tile costs $1.27$ to $2.41$ times the
Clough--Tocher tile of the same region, increasing with
the target; the inverse-distance average costs $1.20$ to $2.01$, very
nearly the same. Relative to each other, the two exceed parity
by only $1.04$ to $1.09$, separable at the three tightest targets and not at
the two loosest. Whatever the roughly two-fold cost at loose targets
represents, it is not the price of restoring spatial detail, because the
transfer that restores none pays it as well. The ratio between the two
bounds the cost of deconvolution rather than isolating it, since the two
chains differ by more than the presence of a deconvolution; read as a
bound, it indicates that this surcharge is small compared with the
surcharge of leaving Clough--Tocher. Isolating it would require a
non-deconvolving transfer of the same smoothness class, which is not
available.

Table~\ref{tab:compactability} isolates that question. Its second column is
each raster's own lossless size; the remaining columns give the compressed
size at each target as a fraction of that same lossless size, so that the
entropy difference between the rasters divides out and what is left is how
much of the content a lossy coder can discard.

\begin{table}[t]
\centering
\caption{Lossless size and compressed fraction, by transfer}
\label{tab:compactability}
\footnotesize
\begin{tabular}{@{}l r rrrrr@{}}
\toprule
 & Lossless & \multicolumn{5}{c}{Fraction of that lossless size} \\
\cmidrule(lr){3-7}
Transfer & (MB) & $0.0001$ & $0.0003$ & $0.001$ & $0.003$ & $0.01$ \\
\midrule
C--T & 1.31 & 0.482 & 0.321 & 0.156 & 0.054 & 0.011 \\
PSF & 1.36 & 0.577 & 0.410 & 0.245 & 0.123 & 0.025 \\
IDW & 1.41 & 0.551 & 0.387 & 0.237 & 0.110 & 0.023 \\
native UTM & 0.37 & 0.819 & 0.584 & 0.349 & 0.157 & 0.034 \\
\bottomrule
\end{tabular}
\vspace{3pt}
\begin{minipage}{\columnwidth}
\footnotesize
Transfers are labelled as in Table~\ref{tab:transfer_cost}. The second column is each raster's own lossless size, which separates the HEALPix rasters by less than a tenth; the rest is the compressed size at each target as a fraction of that same lossless size, so the entropy difference divides out and what remains is how much of the content a lossy coder can discard. Section~\ref{sec:transfer_cost} reads the two surcharges apart.
\end{minipage}
\end{table}


What separates the transfers is not how much each raster carries but how
compactly it can be coded. Their lossless sizes are within $9\%$ of each
other (Table~\ref{tab:compactability}): \SI{1.31}{MB}, \SI{1.36}{MB} and
\SI{1.41}{MB} for Clough--Tocher, the PSF transfer and the inverse-distance
average. Their rate--distortion curves have very different slopes. At the
$10^{-2}$ target the Clough--Tocher tile compresses to $0.8\%$ of its own
lossless size, the deconvolved tile to $1.8\%$, the inverse-distance tile to
$2.0\%$ and the native raster to $3.0\%$ of its own; this is the
oversampling redundancy studied in this paper, which every transfer
inherits. Across regions the surcharge is not explained by how much each
raster carries: correlating each region's own lossless size with the
surcharge it pays gives a Spearman $\rho$ of $-0.46$ for the PSF transfer
($p=0.003$, $n=40$) and $-0.31$ for the inverse-distance average
($p=0.05$, $n=39$), so the regions whose rasters hold the most information
pay the \emph{smallest} surcharge. That is the opposite sign from what an
added-content explanation predicts, and it excludes total information
content as the mechanism. The observations are consistent with energy
compaction being the dominant mechanism, reached by different routes. The deconvolution puts
power back into the small scales directly, $2.80$ times the observation's
in the top octave against $0.89$ for Clough--Tocher, and a coder asked for a
given RMSE has to encode it; the inverse-distance average adds no such
power ($0.92$) and still pays. The cause has not been isolated at the
subband level. What is established is that its output lies ten times
further from an analytic truth than that of Clough--Tocher, and that
content which is neither the scene nor suppressed has to be coded as
content.

The consequence for the headline result is direct, and it is the largest
single sensitivity measured in this study. Under Clough--Tocher the HEALPix
carrier is the cheaper representation at the two loosest targets ($0.93$
and $0.73$ against the native product under the illustrative overhead).
Under either other transfer it is cheaper at none: $2.00$ and $1.72$ for
the PSF transfer, $1.90$ and $1.60$ for the inverse-distance average. A factor of two in the
storage cost of the transfer exceeds the whole margin, so that the
regime-dependent verdict of Section~\ref{sec:rd_main} is a statement about
a carrier populated by this transfer, not about the grid. Stated in final
sizes with no overhead (Table~\ref{tab:headline}): the Clough--Tocher tile
compresses to between $2.06$ and $0.93$ times the native square as the
tolerated error grows, the PSF-aware tile to between $2.62$ and $2.24$
times, and the inverse-distance tile to between $2.49$ and $2.07$ times.
The two latter columns barely move with the tolerated error, which is
another way of saying that their content resists lossy coding, whereas
the content of the Clough--Tocher tile is mostly the smooth redundancy
that a codec is built to remove.

Table~\ref{tab:kernel_ablation} rebuilds the carrier with five kernels from
the same native window, reading each with the same operator, so that any
difference between its rows is the transfer alone. Two of the library's
names mislead and are worth fixing here: \texttt{linear} is the
inverse-distance-weighted average over four neighbours, not a linear
interpolant, and \texttt{bicubic} is Keys' cubic convolution. Its columns
answer four different questions. \emph{Transfer} is the error against the
analytic field, as a percentage of that field's standard deviation.
\emph{Over} counts the cells a transfer places outside the range of the data
it was given, which is how a diverging interpolant shows itself to a check
that a constant field passes. \emph{Floor} is the error of the uncompressed
tile read back on the native pixels, as a fraction of the scene's standard
deviation, with \emph{rel.} its per-region ratio to Clough--Tocher; both
reward a transfer for staying close to the identity on the native lattice,
which is why they are reported and not ranked. \emph{log-PSD} is the RMS
distance in $\log_{10}$ between the tile's radial power spectrum and the
observation's up to the native Nyquist frequency, and \emph{top oct.} the
ratio of their power in the highest octave alone, where a value above one
means the transfer leaves more high-frequency power than the observation
carries. The last three columns are compressed-size ratios to the same
reference at three targets.

\begin{table*}[t]
\centering
\caption{Transfer-kernel ablation}
\label{tab:kernel_ablation}
\footnotesize
\begin{tabular}{@{}lrrrrrr@{}}
\toprule
Kernel & Transfer & Over & Floor & rel. & log-PSD & top oct. \\
\midrule
\texttt{bicubic} & 392.28 & 2.45 & 11.9 & 0.96 & 0.16 & 1.51 \\
\texttt{clough\_tocher} & 1.20 & 0.00 & 10.0 & 1.00 & 0.05 & 0.89 \\
\texttt{linear} & 11.43 & 0.85 & 6.9 & 0.74 & 0.03 & 0.92 \\
\texttt{nearest} & 10.04 & 0.00 & 0.7 & 0.12 & 0.00 & 1.00 \\
\texttt{psf} & 8.16 & 0.19 & 23.6 & 2.55 & 0.44 & 2.80 \\
\bottomrule
\end{tabular}
\vspace{3pt}
\begin{minipage}{\textwidth}
\footnotesize
Medians over the 40 regions of the study (\texttt{bicubic} on 39, \texttt{linear} on 39, the rest on all). Each row rebuilds the carrier tile from the same native window with a different kernel and the same reading operator, so differences between rows are the transfer alone. Two library names mislead: \texttt{linear} is an inverse-distance-weighted average over four neighbours, not a linear interpolant, and \texttt{bicubic} is Keys' cubic convolution, whose signed kernel rings outside the local sample range by design; its row is reported for completeness, not as a finding. \emph{Transfer} is the error against an analytic field known at every position, as a percentage of the field's standard deviation; it is the only column that gives a transfer no credit for copying its input samples. \emph{Over} is the percentage of cells placed outside the range of the data the transfer was given, which detects a diverging interpolant that the constant-field check passes. \emph{Floor} is the RMSE of the uncompressed tile read on the native pixels over the two footprints' intersection, as a fraction of the scene's standard deviation, and \emph{rel.} its per-region ratio to \texttt{clough\_tocher}; both reward a transfer for staying close to the identity on the native lattice and are reported for that reason, not as a ranking. \emph{log-PSD} is the RMS distance in $\log_{10}$ between the projected tile's radial power spectrum and the observation's up to the native Nyquist, and \emph{top oct.} the ratio of their power in the highest octave alone: above 1 the transfer leaves more high-frequency power than the observation carries.
\end{minipage}
\end{table*}


Equal bytes are not equal value, and this is where the two surcharges
differ. The accuracy column of Table~\ref{tab:kernel_ablation} scores each
transfer against an analytic field known at every position, the only
measure in that table that gives a transfer no credit for copying its input
samples. Among the transfers that aim to reproduce that field,
Clough--Tocher reaches $1.20\%$ of its standard deviation, against
$10.04\%$ for nearest and $11.48\%$ for the inverse-distance average, on the
39 regions the window audit of Appendix~\ref{app:window_radiometry} admits
(38 for the inverse-distance average and for bicubic). The
bicubic row is reported for completeness: the signed kernel of Keys rings
outside the local sample range by design, and its divergence at this
sampling ratio is a property of that choice rather than of the carrier
grid. The $8.15\%$ of the PSF transfer is not a ranking entry, since its
estimand is the latent field and departing from the observed one is its
purpose. What the column shows is that the additional bytes of the
inverse-distance average buy transfer error, tenfold.

The additional bytes of the deconvolved carrier are of the opposite kind,
and this is established by the companion study rather than by anything
measured here. On a known pre-degradation texture the matched
reconstruction recovers $15.4\%$ to $19.6\%$ of the imposed blur error, and
on real reduced-resolution Sentinel-2 imagery it has the smallest RMSE in
all 40 regions, improving on the best of nearest, bilinear, bicubic and
Richardson--Lucy by $11.0\%$ to $17.9\%$, with the restored power at low and
intermediate frequencies rather than concentrated at the Nyquist
frequency~\cite{delouis2026resample}. This is a statement about the method
under the conditions of that study; whether the particular tiles whose
bytes are counted here realise it is not established, and our own harness
does not reproduce the magnitude of the recovery
(Appendix~\ref{app:limitations}). With this reservation, the two-fold
surcharge has two different meanings: on the deconvolved carrier it is the
price of small-scale content placed there by a validated method, on the
inverse-distance carrier the price of content that the ablation of this
paper shows to lie ten times further from the truth than the transfer it
replaced. Clough--Tocher remains both the most accurate of the reproducing
transfers and the most compressible, and the two properties have a common
cause: a transfer that represents a band-limited field faithfully does not
invent fine-scale content, and content that it does not invent is content
that the codec is not asked to pay for. The scope of the result is
therefore narrower than a grid comparison and more useful than a kernel
preference. A HEALPix level-20 carrier becomes competitive at aggressive
compression when it is populated by a transfer in the smoothness class of
Clough--Tocher; leaving that class costs about a factor of two at the
loosest targets and removes the advantage; and this cost is a property of
how compactly the output of the transfer can be coded, not of how much
detail it restores.
